% il existe plusieures classes de documents
% pour des documents plus longs, vous pouvez utiliser
% book ou report
\documentclass[a4paper]{article} 
\usepackage{graphicx}
\usepackage{biblatex} %Imports biblatex package
\usepackage{subfig}
\usepackage{parskip}
\usepackage{multirow}
\usepackage{multirow}
\usepackage[most]{tcolorbox}
\usepackage{xcolor}
\addbibresource{stg3a.bib} %Import the bibliography file
\usepackage{hyperref}

% vous pouvez changer les paramètres : voici les options dispoinibles :
% - a4paper
% - fancysections
% - notitlepage ou titlepage
% - onside ou twoside selon si vous voulez l'imprimer en recto-verso ou en recto
% - sectionmark
% - chaptermark (pour les 
% - pagenumber
% - enmanquedinspiration
% en cas de doutes, pas de doutes, la documentation est sur :
%  https://gitlab.binets.fr/typographix/polytechnique/-/blob/master/source/polytechnique.pdf



\usepackage[a4paper,  fancysections,  titlepage]{polytechnique}
\usepackage[french]{babel}
\usepackage[T1]{fontenc}
\usepackage{blindtext}
\usepackage[hidelinks]{hyperref}

\usepackage{listings}
\usepackage{xcolor}
\lstset{
  language=Python,
  basicstyle=\ttfamily\small,
  breaklines=true,
  showstringspaces = false,
  keywordstyle=\color{blue},
  stringstyle=\color{red},
  commentstyle=\color{brown},
  morecomment=[l][\color{magenta}]{\#},
  inputencoding=utf8,
  extendedchars=true,
  literate={à}{{\`a}}1 {ç}{{\c{c}}}1 {é}{{\'e}}1 {è}{{\`e}}1 {ê}{{\^e}}1 {ë}{{\"e}}1 {î}{{\^i}}1 {ï}{{\"i}}1 {ô}{{\^o}}1 {ö}{{\"o}}1 {ù}{{\`u}}1 {û}{{\^u}}1 {ü}{{\"u}}1 {ÿ}{{\"y}}1 {À}{{\`A}}1 {Ç}{{\c{C}}}1 {É}{{\'E}}1 {È}{{\`E}}1 {Ê}{{\^E}}1 {Ë}{{\"E}}1 {Î}{{\^I}}1 {Ï}{{\"I}}1 {Ô}{{\^O}}1 {Ö}{{\"O}}1 {Ù}{{\`U}}1 {Û}{{\^U}}1 {Ü}{{\"U}}1 {Ÿ}{{\"Y}}1 {\_}{\textunderscore}{1}
}

% nous avons défini deux commandes :
\newcommand{\code}[1]{%
    \mbox{\ttfamily%
        \detokenize{#1}%
    }%
}

\newcommand{\resultat}[1]{%
    \quad \rightsquigarrow \quad #1%
}

\author{Auteur}
\date{\today}
\title{Titre}
\subtitle{Sous-titre}%
% pour changer de logo, ajoutez l'image dans un fomat PDF
% ou image en la glissant à droite et remplacez typographix
% par le nom de l'image, si vous ne voulez pas de logo, 
% supprimze la ligne. 

\begin{figure*}[b]
    \includegraphics[width=0.4\linewidth]{logoaividence.jpg}
    \label{fig:enter-label}
\end{figure}


\fontsize{40pt}{48pt}

\begin{document}
	\maketitle 
	
	\tableofcontents
\newpage

\setlength{\parindent}{2em}


\addcontentsline{toc}{section}{Résumé}
\section*{Résumé}
\setlength{\parindent}{2em}
(paragraphe)


\newpage
\addcontentsline{toc}{section}{Executive summary}
\section*{Executive summary}
(paragraphe)


\newpage
\addcontentsline{toc}{section}{Introduction}
\section*{Introduction}
(paragraphe)

\newpage

\addcontentsline{toc}{section}{Notations}
\section*{Notations}
(paragraphe)

\newpage

\section{Titre}

\subsection{Titre}
(paragraphe)


\subsection{Titre}
(paragraphe)

\subsubsection{Titre}
(paragraphe)

\begin{center}
    \textbf{Analyse PESTEL}
\end{center}

\begin{itemize}

\item{$\bullet$  \textbf{Politique :} }
(paragraphe)
\item{$\bullet$  \textbf{Economie :} }
(paragraphe)
\item{$\bullet$  \textbf{Sociétal et Culturel :} }
(paragraphe)
\item{$\bullet$ \textbf{Technologique: }} 
(paragraphe)
\item{$\bullet$  \textbf{Environnemental :} } 
(paragraphe)
\item{$\bullet$  \textbf{Légal : }} 
(paragraphe)

(paragraphe)
\end{itemize}

\subsubsection{Titre}
\begin{itemize}
\item{$\bullet$ \textbf{Risque d'entrée de nouveaux acteurs :} } 
(paragraphe)
\item{$\bullet$ \textbf{Pouvoir de négociation des fournisseurs :} } 
(paragraphe)
\item{$\bullet$ \textbf{Pouvoir de négociation des clients :} } 
(paragraphe)
\item{$\bullet$ R\textbf{isque de produits alternatifs :} }
(paragraphe)
\item{$\bullet$ \textbf{Concurrence entre les acteurs établis :} }
(paragraphe)

(paragraphe)
\end{itemize}

\subsubsection{Titre}
(paragraphe)
\begin{center}
\textbf{Évaluation SWOT}
\end{center}

\begin{itemize}
\item \textbf{$\bullet$ Atouts :} 
 \textit{Compétence en Explicabilité :} (paragraphe) 
 \textit{Conseil Scientifique Prestigieux :} (paragraphe) 
 \textit{Approche Unique :} (paragraphe) 
 \textit{Incubation au sein de Paris-Saclay :} (paragraphe) 
 \textit{Alliances Stratégiques :} (paragraphe) 

\item \textbf{$\bullet$ Lacunes :} 
 \textit{Équipe Restreinte :} (paragraphe) 
 \textit{Dépendance aux Stagiaires :} (paragraphe) 

\item \textbf{$\bullet$ Opportunités :} 
 \textit{Demande Accrue pour l’Explicabilité :} (paragraphe) 
 \textit{Évolution Réglementaire :} (paragraphe) 

\item \textbf{$\bullet$ Risques :} 
 \textit{Accroissement de la Concurrence :} (paragraphe)
 \textit{Contraintes Régionales :} (paragraphe) 
 \textit{Dépendance aux Fournisseurs et aux Ressources :} (paragraphe)
\end{itemize}
(paragraphe)

\begin{figure}[h]
    \centering
    \includegraphics[width=0.8\linewidth]{Compétence en Explicabilité Conseil scientifique prestigieux Quelles sont les plus grandes qualités de l’entreprise .jpg}
    \caption{Titre}
    \label{fig:enter-label}
\end{figure}


\subsection{Titre}

\subsubsection{Titre}

(paragraphe)

\textbf{$\bullet$ Organisation hiérarchique :} (paragraphe)
\begin{itemize}
    \item \textbf{David Cortés}, (paragraphe)
    \item \textbf{Laurent}, (paragraphe)
\end{itemize}
(paragraphe)

\textbf{$\bullet$ Conseil scientifique :} (paragraphe)
\begin{itemize}
    \item \textbf{Pierre Marquis}, (paragraphe)
    \item \textbf{Arthur Charpentier}, (paragraphe)
    \item \textbf{Stéphane Clémençon}, (paragraphe)
\end{itemize}
(paragraphe)

\textbf{$\bullet$ Mécanismes de Coordination selon MINTZBERG :} (paragraphe)
(paragraphe)

\subsubsection{Titre}

\textbf{$\bullet$ La culture d'entreprise chez AI-VIDENCE :} (paragraphe)
\textbf{$\bullet$ Signes :} (paragraphe)
\textbf{$\bullet$ Symboles et Rituels :} (paragraphe)
\textbf{$\bullet$ Vision Collective :} (paragraphe)

\subsubsection{Titre}

\textbf{$\bullet$ Motivations et incitations :}
(paragraphe)
— \textit{$\bullet$ Mission Porteuse de Sens :} (paragraphe) 
— \textit{$\bullet$ Récompenses liées à la Performance :} (paragraphe) 
— \textit{$\bullet$ Souplesse et Valorisation :} (paragraphe) 
— \textit{$\bullet$ Cohésion par les Activités Sociales :} (paragraphe)
(paragraphe)

\textbf{$\bullet$ Répartition du pouvoir :}
(paragraphe)

\newpage

\section{Titre}

\subsection{Titre}
(paragraphe)

\subsection{Titre}
\subsubsection{Titre}
(paragraphe)

\begin{figure}[h]
    \centering
    \includegraphics[width=0.7\linewidth]{Capture d'écran 2024-09-22 160601.png}
    \caption{Titre}
    \label{fig:enter-label}
\end{figure}

\subsubsection{Titre}
(paragraphe)

\begin{itemize}
    \item \textbf{$\bullet$ Découverte du cadre de la mission :}
    (paragraphe)
    \item \textbf{$\bullet$ Échanges sur les attentes et objectifs :}
    (paragraphe)
    \item \textbf{$\bullet$ Analyse du produit AntakIA :}
    (paragraphe)
    \item \textbf{$\bullet$ Examen des projets antérieurs :} 
    (paragraphe)
    \item \textbf{$\bullet$ Participation à des événements extérieurs :}
    (paragraphe)
\end{itemize}

\paragraph{Titre}
(paragraphe)

\begin{figure}[h]
    \centering
    \includegraphics[width=0.7\linewidth]{Capture d’écran 2024-09-22 160656.png}
    \caption{Titre}
    \label{fig:enter-label}
\end{figure}

\subsubsection{Titre}
\setlength{\parindent}{2em}
(paragraphe)
\begin{itemize}
  \item \textbf{$\bullet$ Examen de l'approche du projet HierarchicalCausalModels (HCM) : }
  (paragraphe)
  \item \textbf{$\bullet$ Exploration de la littérature scientifique :} 
  (paragraphe)
  \item \textbf{$\bullet$ Mise en œuvre de la modélisation causale hiérarchique : }
  (paragraphe)
  \item \textbf{$\bullet$ Intégration dans CausaLink : }
  (paragraphe)
\end{itemize}

\paragraph{\textbf{Titre}} 
(paragraphe)

\subsubsection{Titre}
(paragraphe)
\begin{itemize}
    \item \textbf{$\bullet$ Documentation détaillée :}
    (paragraphe)
    \item \textbf{$\bullet$ Suivi des projets :} 
    (paragraphe)
    \item \textbf{$\bullet$ Tests et validation :}
    (paragraphe)
\end{itemize}

\paragraph{\textbf{Titre}} (paragraphe)

\newpage
\section{Titre}
(paragraphe)

\subsection{Titre}
(paragraphe)
\subsection{Titre}
(paragraphe)
\subsection{Titre}
(paragraphe)
\subsection{Titre}
(paragraphe)
\subsection{Titre}
(paragraphe)

\newpage

\addcontentsline{toc}{section}{Conclusion}
\section*{Conclusion}
\setlength{\parindent}{2em}
(paragraphe)

\newpage

\section*{Annexes}

\begin{figure}[h]
    \centering
    \includegraphics[width=0.8\linewidth]{001504621_896x598_c.png}
    \caption{Titre}
    \label{fig:enter-label}
\end{figure}


\newpage

\appendix
\section{Expériences numériques}
(paragraphe)

\printbibliography

\end{document}
